\scanpage{223}
\mspage{78}

entre l'\unsure{ens.} des \unsure{élément|éléments} clos de $P$, et l'ens des \unsure{éléments} clos de $Q$

\diagram{\begin{tikzcd}[ampersand replacement=\&,column sep=large]
\text{(5.22.)} \& \unsure{p|P}^{\unsure{\sharp}}
 \arrow[r,shift left,"{\substack{\overline{\beta}_0\\\sim}}"] \&
 Q^\flat\arrow[l,shift left,"\overline{\alpha}_0","\sim"'] .
\end{tikzcd}}

Mais en pratique, les applications $\alpha_0,\beta_0$ ne sont pas obtenues directement, mais on dispose de la relation $R$, satisfaisant la condition $(D_0)$, dont on déduit $\alpha_0,\beta_0$. Dans tous les cas auxquels je \deleted{\unsure{connais}} \inserted{pense} ici, ou bien $P,Q$ \unsure{sont} de la forme $\unsure{\mathfrak P|\mathcal P}(P_0),\mathfrak P(Q_0)$, et $R$ se déduit d'une relation $R_0\subset P_0\times Q_0$ comme ci-dessus (cas de loin le plus fréquent, il me semble), ou bien $P$ et $Q$ apparaissent comme sous-ensembles \inserted{convenables} d'ensembles $\mathfrak P(P_0),\mathfrak P(Q_0)$, munis de l'ordre induit, et la relation $R$ sur $P\times Q$ s'obtient d'une façon plus complexe.

\uline{Exemples}

1) Relation d'\uline{orthogonalité} entre \deleted{parties} de deux modules sur un anneau, \deleted{\unsure{munis d'un}} d'un accouplement de dualité. \inserted{Dans le cas de \unsure{vectoriels|modules} de dimension finie et d'une vraie dualité, les parties closes sont les sous-esp. vectoriels.}

% margin note written sideways
\margin{\deleted{\unsure{dans, ce sont}} dans \unsure{le|ce} cas \unsure{infini|en fait}, \unsure{se|ce} sont les sous-espaces vectoriels fermés pour la ``\unsure{topologie|topol.} \unsure{faible}'' définie par l'accouplement.}

2) Relations d'\uline{orthogonalité} entre parties de deux groupes topologiques \unsure{comm.} loc. \unsure{cpcts|cpacts} duaux l'un de l'autre, au sens de \unsure{Pontragin|Pontriagin|Pontrjagin|Pontraguin}.
% archive sheet number: 223
